\documentclass[11pt]{article}

\usepackage[utf8]{inputenc}
\usepackage[T1]{fontenc}
\usepackage{lmodern}
\usepackage{geometry}
\usepackage{amsmath,amssymb,amsfonts,amsthm,mathtools}
\usepackage{array}
\usepackage{enumitem}
\usepackage{microtype}
\usepackage{hyperref}

\geometry{margin=1in}
\hypersetup{colorlinks=true, linkcolor=black, urlcolor=blue, citecolor=black}
\setlist[enumerate]{topsep=0.35em,itemsep=0.3em}
\setlist[itemize]{topsep=0.3em,itemsep=0.25em}

\newcommand{\Obj}{\mathsf{Obj}_{0}^{\mathrm{iso}}}
\newcommand{\Inv}{\mathsf{Inv}_{0}^{\mathrm{iso}}}
\newcommand{\Cmp}{\mathsf{Cmp}_{0}^{\mathrm{iso}}}
\newcommand{\EqSrc}{\mathsf{Eq}_{\mathrm{src}}}
\newcommand{\EqIso}{\mathsf{Eq}_{\mathrm{src}}^{\mathrm{iso}}}
\newcommand{\GenH}{\mathsf{Gen}_{H}}
\newcommand{\RetainH}{\mathsf{Retain}_{H}}
\newcommand{\NoTargetH}{\mathsf{NoTarget}_{H}}
\newcommand{\NoTargetSyn}{\mathsf{NoTarget}_{H}^{\mathrm{syn}}}
\newcommand{\SrcTok}{\mathsf{Tok}_{0}}
\newcommand{\ForbidT}{\mathsf{Forbid}_{T}}
\newcommand{\ProxyT}{\mathsf{Proxy}_{T}}
\newcommand{\TraceH}{\mathsf{Trace}_{H}}
\newcommand{\Annot}{\mathsf{Ann}}
\newcommand{\Influence}{\mathsf{Infl}}
\newcommand{\FreezeH}{\mathsf{Freeze}_{H}}
\newcommand{\JudgeSyn}{\mathsf{J}_{\mathrm{syn}}}
\newcommand{\Pass}{\mathsf{Pass}}
\newcommand{\Block}{\mathsf{Block}}
\newcommand{\Iso}{\mathsf{Iso}_{0}}

\newtheorem{definitionblock}{Definition}
\newtheorem{datum}{Datum}
\newtheorem{ruleblock}{Rule}
\newtheorem{finding}{Finding}
\newtheorem{verdict}{Local verdict}
\newtheorem{nextburden}{Next burden}

\title{\textbf{Localization Source-Basis Axiom Selector-Domain EqSrc Source-Relabeling Isomorphism Datum}\\[0.35em]
\large Bounded Ontology Formalizer Evidence-Supply Packet for the Manifest-Bound \(\NoTargetSyn\) Trace}
\author{Ontology Formalizer AgentJob AJ-RT-20260613-026-001}
\date{June 14, 2026}

\begin{document}

\maketitle

\begin{abstract}
This draft/control artifact supplies one bounded source-only
\(\EqSrc\)-datum at the weakest available level: source-relabeling
isomorphism for the manifest-bound \(\NoTargetSyn\) trace and its
structure-preserving relabeling. The datum provides a local object domain,
witness pair, invariant ledger, comparison rule, annotation obligations,
influence-edge coverage, and negative controls. It supports only a local
\(\EqIso\) judgment for the original frozen trace/manifest and a
source-neutral relabeled copy. It does not discharge general \(\EqSrc\),
\(\RetainH\), or \(\GenH\); it does not prove semantic \(\NoTargetH\);
and it does not authorize candidate reconstruction, canonical ontology
edit, benchmark promotion, Gate Chair review, or completed-derivation
language.
\end{abstract}

\section{Scope}

This artifact is the role output for task \texttt{RT-20260613-026}. It
answers the narrow evidence-supply question left by
\texttt{handoff-0042}: can any source-only \(\EqSrc\) datum be registered
without target-channel or proxy selection?

The answer is positive only in a local and syntactic sense. The available
tracked evidence supports a source-relabeling isomorphism datum over the
registered \(\NoTargetSyn\) trace and manifest. It does not support a
general equivalence relation over generated ledgers, candidate families, or
empirical behaviors.

\section{Source Objects}

\begin{definitionblock}[Registered base object]
Let
\[
    X_{0}=(\TraceH,\FreezeH^{\mathrm{reg}}_{\mathrm{RT21}})
\]
where
\[
    \FreezeH^{\mathrm{reg}}_{\mathrm{RT21}} =
    (\SrcTok^{\mathrm{RT21}},\ForbidT^{\mathrm{RT21}},
    \ProxyT^{\mathrm{RT21}},\Annot^{\mathrm{RT21}},
    \Influence^{\mathrm{RT21}}).
\]
The object \(X_{0}\) is the manifest-bound source-syntax object used in
the \(\NoTargetSyn\) theorem retry and stress test.
\end{definitionblock}

\begin{definitionblock}[Source-neutral relabeling object]
Let \(v_{\mathrm{iso}}\) be a bijective relabeling of ST, PT, AN, and IE
identifiers that preserves token class, annotation content, forbidden
channel status, proxy class, and influence-edge incidence. Define
\[
    X_{1}=v_{\mathrm{iso}}(X_{0}).
\]
The object \(X_{1}\) is not a new physical candidate and not a canonical
ontology object. It is a source-syntax relabeling witness.
\end{definitionblock}

\begin{datum}[Local object domain]
The object domain for this datum is
\[
    \Obj=\{X_{0},X_{1},X_{\mathrm{ann}},X_{\mathrm{proxy}}\},
\]
where \(X_{\mathrm{ann}}\) is an annotation-weakening variant and
\(X_{\mathrm{proxy}}\) is an unmapped-proxy variant. The last two objects
are negative controls. They are included only to test that the comparison
rule can return \(\Block\) for source-side reasons.
\end{datum}

\section{Witness Pair}

\begin{datum}[Positive witness pair]
The positive witness pair is
\[
    (X_{0},X_{1})\in \Obj\times\Obj.
\]
Its intended local judgment is
\[
    \EqIso(X_{0},X_{1}).
\]
The pair is selected because `RT-20260613-024` records that the
\(\NoTargetSyn\) source-syntax pass survives structure-preserving
source-token relabeling. The selection criterion is manifest structure,
not target metric behavior, target proper time, target locality, target
observer protocol, target-favorable rank, benchmark-success evidence, or
candidate performance.
\end{datum}

\begin{datum}[Negative controls]
The negative-control pairs are
\[
    (X_{0},X_{\mathrm{ann}}),\qquad
    (X_{0},X_{\mathrm{proxy}}).
\]
The first removes required path, hash, rule-scope, dependency, exclusion,
proxy, influence, or manifest-use annotation content. The second introduces
or uses an unmapped proxy class. The intended result for both pairs is
\(\Block\), not \(\EqIso\).
\end{datum}

\section{Invariant Ledger}

\begin{datum}[Source-only invariants]
The comparison rule may use only the following source-side invariants:
\begin{enumerate}
    \item \textbf{Token-class incidence}: ST, PT, AN, and IE identifiers
    must map bijectively to identifiers of the same class.
    \item \textbf{Forbidden-channel orientation}: FT-001--FT-010 must
    remain blocker channels rather than admitted source primitives.
    \item \textbf{Proxy-class coverage}: PT-001--PT-010 must remain mapped
    as suspect proxy classes, including PT-002 for source-equivalence
    proxy risk and PT-007 for notation-similarity risk.
    \item \textbf{Annotation content}: AN-001--AN-010 must preserve source
    path, object identifier, hash, timing, rule scope, dependency,
    exclusion, proxy classification, influence-edge content, and manifest
    use limits.
    \item \textbf{Influence-edge incidence}: IE-001--IE-014 must preserve
    source and destination token classes and blocker state.
    \item \textbf{No closed-world inference}: omitted tokens, proxies, or
    influence edges may not be treated as absent by default.
    \item \textbf{Local judgment shape}: the source-syntax adjudicator
    shape must remain a manifest-bound syntactic \(\Pass\) or \(\Block\)
    rule, not semantic target comparison.
\end{enumerate}
\end{datum}

These invariants are source-side because they inspect only registered
token class, manifest membership, annotation content, source hashes, edge
incidence, blocker state, and rule scope. They do not inspect target
outcomes or candidate performance.

\section{Comparison Rule}

\begin{ruleblock}[Local relabeling comparison]
For \(X,Y\in\Obj\), define \(\Cmp(X,Y;\Inv)\) as follows:
\begin{enumerate}
    \item Return \(\EqIso(X,Y)\) if there exists a bijection \(\rho\) on
    the ST, PT, AN, and IE identifiers such that every invariant in
    \(\Inv\) is preserved under \(\rho\).
    \item Return \(\Block\) if either object uses an unlisted token,
    admits a forbidden target channel as a source primitive, uses an
    unmapped proxy, weakens a required annotation, changes influence-edge
    incidence, infers edge absence from silence, or requires semantic
    target comparison.
    \item Return \(\neg\EqIso(X,Y)\) only when both objects are well-formed
    source objects under the manifest universe but no invariant-preserving
    bijection exists.
\end{enumerate}
\end{ruleblock}

\begin{finding}[Mechanically checkable status]
The positive case is mechanically checkable from the registered
\(\NoTargetSyn\) freeze manifest and stress-test statement because it asks
only whether source token classes, annotations, forbidden-channel
orientation, proxy coverage, and influence incidence are preserved by a
renaming. The negative controls are also mechanically checkable: missing
annotation content and unmapped proxy use force \(\Block\).
\end{finding}

\section{Annotation and Influence Receipts}

\begin{datum}[Annotation receipts]
The local datum inherits the following required annotation surfaces:
\[
\begin{array}{ll}
AN\text{-}001 & \text{source path object id hash and admission rule},\\
AN\text{-}002 & \text{comparison rule and proof that no target proxy selects the comparison},\\
AN\text{-}008 & \text{proxy classification affected decision class and blocker state},\\
AN\text{-}009 & \text{source and destination token declared rule and blocker state},\\
AN\text{-}010 & \text{manifest-use limit preserving theorem-discharge blocks}.
\end{array}
\]
For this datum, AN-002 is discharged only locally: the comparison rule is
selected by source-relabeling structure and returns \(\Block\) for proxy or
annotation weakening.
\end{datum}

\begin{datum}[Influence receipts]
The local comparison explicitly uses IE-006 for PT-002 to ST-002,
IE-011 for PT-007 to ST-002 and ST-005, IE-012 for validation-passage
proxy surfaces, IE-013 for human-authorization proxy surfaces, and IE-014
for generated-derivative visibility proxy surfaces. These are guard edges,
not positive equivalence evidence.
\end{datum}

\section{Local Result}

\begin{verdict}[Local EqSrc datum supplied]
The packet supplies the local source-only datum
\[
    \mathcal{E}_{0}^{\mathrm{iso}} =
    (\Obj,\Inv,\Cmp,\Annot,\Influence)
\]
for the witness pair \((X_{0},X_{1})\). Within this bounded datum,
\[
    \Cmp(X_{0},X_{1};\Inv)=\EqIso(X_{0},X_{1}).
\]
\end{verdict}

\begin{verdict}[Negative controls block]
For the annotation-weakening and unmapped-proxy variants,
\[
    \Cmp(X_{0},X_{\mathrm{ann}};\Inv)=\Block,\qquad
    \Cmp(X_{0},X_{\mathrm{proxy}};\Inv)=\Block.
\]
The block reason is source-side: required annotation content or mapped
proxy coverage is absent.
\end{verdict}

\begin{verdict}[No general EqSrc discharge follows]
This result does not prove that \(\EqSrc\) is a global equivalence relation
over generated ledgers or candidate families. It does not establish
reflexivity, symmetry, transitivity, quotient preservation, adverse-output
retention, complete negative-control retention, or generator adequacy
outside the local source-relabeling pair.
\end{verdict}

\section{Boundary Ledger}

\begin{center}
\begin{tabular}{p{0.25\linewidth}p{0.43\linewidth}p{0.18\linewidth}}
\textbf{Surface} & \textbf{Result} & \textbf{Status} \\
\hline
\(\Obj\) & local domain with one positive relabeling pair and two negative controls & supplied \\
\(\Inv\) & source-only token class annotation proxy and influence invariants & supplied \\
\(\Cmp\) & mechanical relabeling-isomorphism rule with Block conditions & supplied \\
AN-002 & comparison rule and proxy-exclusion proof only for local pair & local only \\
IE-006 & PT-002 to ST-002 guard edge & guard only \\
\(\EqIso(X_{0},X_{1})\) & local source-relabeling judgment & draft/control evidence \\
General \(\EqSrc\) & equivalence relation over full source object families & not discharged \\
\(\RetainH,\GenH\) & not attempted & blocked \\
Candidate reconstruction & not authorized & blocked \\
Gate Chair review & not authorized and human-gated & blocked \\
\end{tabular}
\end{center}

\section{Next Burden}

\begin{nextburden}[Audit before broader equivalence use]
The logical next role is Smuggling Auditor. The audit should test whether
this local \(\EqIso\) datum launders target success through notation,
manifest choice, proxy selection, annotation preservation, influence-edge
selection, or authority labels. Until that audit passes, and until a later
bounded task proves broader equivalence-relation properties and retention
coverage, \(\EqSrc\), \(\RetainH\), \(\GenH\), candidate reconstruction,
canonical ontology edit, benchmark promotion, Gate Chair review, and
completed-derivation language remain blocked.
\end{nextburden}

\section*{References}

Ricciardi, A. S. (2026, June 14). \emph{Localization source-basis axiom
selector-domain EqSrc source-equivalence obligation packet or controlled
pause: Ontology Formalizer review after the stress-tested
\(\NoTargetSyn\) source-syntax guard} [Unpublished manuscript]. The
Aether-Flow Interpretation of Relativity Research Project,
\texttt{research\_control/tasks/RT-20260613-025/artifacts/41\_LOCALIZATION\_SOURCE\_BASIS\_AXIOM\_SELECTOR\_DOMAIN\_EQSRC\_SOURCE\_EQUIVALENCE\_OBLIGATION\_PACKET\_OR\_CONTROLLED\_PAUSE.tex}.

Ricciardi, A. S. (2026, June 14). \emph{Localization source-basis axiom
selector-domain NoTarget source-syntax benchmark-independence and
proxy-variation stress test: Refuter review of the manifest-bound
\(\NoTargetSyn\) source-syntax pass} [Unpublished manuscript]. The
Aether-Flow Interpretation of Relativity Research Project,
\texttt{research\_control/tasks/RT-20260613-024/artifacts/40\_LOCALIZATION\_SOURCE\_BASIS\_AXIOM\_SELECTOR\_DOMAIN\_NO\_TARGET\_SOURCE\_SYNTAX\_BENCHMARK\_INDEPENDENCE\_PROXY\_VARIATION\_STRESS\_TEST.tex}.

Ricciardi, A. S. (2026, June 13). \emph{Localization source-basis axiom
selector-domain NoTarget source-syntax theorem retry: Manifest-bound
\(\NoTargetSyn\) adjudication after source-registered freeze audit}
[Unpublished manuscript]. The Aether-Flow Interpretation of Relativity
Research Project,
\texttt{research\_control/tasks/RT-20260613-023/artifacts/39\_LOCALIZATION\_SOURCE\_BASIS\_AXIOM\_SELECTOR\_DOMAIN\_NO\_TARGET\_SOURCE\_SYNTAX\_THEOREM\_RETRY.tex}.

\end{document}
